% Capítulo 2 (Introdução) — Slide — Objetivo da pesquisa
\begin{frame}{Objetivo da pesquisa}
  \begin{center}
    \textcolor{darkblue}{\textbf{Objetivo geral}}\\[0.4em]
    \large Formular e avaliar o AIPE como um \textit{framework} metodológico
    para \highlight{seleção contextual de políticas de reposição} em séries
    loja-produto de redes varejistas heterogêneas.
  \end{center}
  \vspace{0.6em}
  \textcolor{darkblue}{\textbf{Objetivos específicos}}
  \begin{enumerate}\setlength\itemsep{1pt}
    \item Caracterizar séries loja-produto a partir de variáveis operacionais
      de demanda.
    \item Avaliar um portfólio de políticas de reposição em ambiente de
      simulação controlado.
    \item Gerar rótulos de política viável de menor CTI sob restrição de
      nível de serviço.
    \item Estruturar a base metodológica para treinamento e validação futura
      do PSE.
  \end{enumerate}
  \note{
    O objetivo geral traduz a pergunta de pesquisa em um programa de
    trabalho: formular e avaliar o AIPE como framework de seleção
    contextual. Os quatro objetivos específicos seguem a ordem em que o
    pipeline metodológico é executado: primeiro caracterizamos as séries
    loja-produto; depois avaliamos o portfólio de políticas em simulação
    controlada; em seguida geramos os rótulos de política viável de menor
    CTI, sempre sob a restrição de nível de serviço; e por fim estruturamos
    a base metodológica -- características mais rótulo -- que sustentará o
    treinamento e a validação futura do PSE, ainda não realizados nesta
    etapa.
  }
\end{frame}
